\uuid{6Q1s}
\titre{Probabilité d'une queue droite pour une loi normale}
\auteur{Maxime Nguyen}
\organisation{AMSCC}
\datecreate{2026-11-27}
\niveau{L2}
\module{Probabilité}
\chapitre{Probabilité continue}
\sousChapitre{Loi normale}
\difficulte{1}
\theme{loi normale, standardisation, symétrie, complémentaire}

\contenu{
\begin{question}
On note $\mathcal{N}(\mu,\sigma)$ la loi normale d'espérance $\mu$ et d'écart-type $\sigma$, et $\Phi$ la fonction de répartition de $\mathcal{N}(0,1)$. Soit $X \sim \mathcal{N}(50, 10)$. Quelles égalités sont vraies ?
\begin{answers}
\bad{$\mathbb{P}(X > 65) = \Phi(1{,}5)$}
\good{$\mathbb{P}(X > 65) = 1 - \Phi(1{,}5)$}
\bad{$\mathbb{P}(X > 65) = 1 - \Phi(15)$}
\good{$\mathbb{P}(X > 65) = \mathbb{P}(X < 35)$}
\good{$\mathbb{P}(X > 65) = \Phi(-1{,}5)$}
\end{answers}
\begin{explanations}
On standardise : $\frac{65 - 50}{10} = 1{,}5$ et
$\mathbb{P}(X > 65) = \mathbb{P}(Z > 1{,}5) = 1 - \Phi(1{,}5) = \Phi(-1{,}5)$
pour $Z \sim \mathcal{N}(0,1)$.

A donne $\mathbb{P}(X \leq 65)$ : on a oublié le complémentaire.

C : on a centré sans diviser par l'écart-type.

D : la densité de $X$ est symétrique par rapport à $\mu = 50$, et $65$ et $35$ sont à la même distance $15$ de $50$.

Contre-sens visés : oublier le complémentaire pour un événement « $>$ » ;
standardiser à moitié ; ne voir la symétrie qu'autour de $0$.
\end{explanations}
\end{question}
}
